\chapter{晶体、玻璃与新材料}

\begin{kaichang}
找一支铅笔，在纸上轻轻划一道。黑线出现了。你可能想不到：这道一蹭就脏手的黑痕，和珠宝店里那颗闪闪发光、号称“自然界最硬”的钻石，居然是同一种元素——碳——组成的，而且纯度都不低。

先别急着翻答案，猜一猜：为什么同样都是碳原子，一个软到能在纸上掉渣、用手指都能捻碎，另一个却硬到能划开玻璃、切开石头？再猜第二个问题：铅笔芯里真的有“铅”吗？

还有你家那扇玻璃窗。它透明、脆、敲一下就碎。可有种流传很广的说法：玻璃其实不是固体，而是一种流动得极其缓慢的液体——“你看欧洲老教堂的窗户，底下比上头厚，就是几百年流下来的”。这是真的吗？

这一章结束时，三个问题都会有答案。而答案藏在同一个地方：原子是\textbf{怎么排队}的。
\end{kaichang}

\section{看得见的秩序}

抓一把食盐倒在黑纸上，用放大镜仔细看：一粒一粒全是小小的立方体，边角整齐得像用尺子切出来的。不信你砸碎一粒（隔着纸用勺背压），碎片还是更小的立方体。

这不只是盐的癖好。大块冰糖上能看到平整光滑的晶面；雪花千差万别，却永远是六瓣；云母能像撕便利贴一样剥成透明的薄片；水晶会自己长成六棱柱。天然矿物这种“自己长出规则外形”的本事，古人早就注意到了，但几千年都没说出个所以然。你手腕上若戴着石英表，表里那片人工切割的小石英晶体，此刻正像音叉一样以极稳的频率振动，替你数着秒——晶体的规则结构，就是它守时的本钱。

这些晶体还有一个共同点：\textbf{熔点固定}。冰在 $0\,^{\circ}\mathrm{C}$ 准时化成水，食盐要烧到 $801\,^{\circ}\mathrm{C}$ 才熔化，熔化过程里温度死死钉在那条线上不动。

再对比另一族物质：玻璃、石蜡、松香、塑料。给它们加热，没有“到哪一度突然熔化”的界线——先变软，再变稠，最后才流动。敲碎一块玻璃，断面是弯弯曲曲的贝壳状曲面，找不到一个平整的平面。蜡块干脆可以被你捏变形。

生活里还有介于两者之间的例子：巧克力。一块好巧克力掰开时有清脆的响声，断面光滑发亮，含在嘴里刚好在体温附近熔化；可放久了表面会起一层白霜，口感变沙。那不是发霉（还能吃），是里面的脂肪悄悄换了一种“排队方式”——记住这个线索，本章后面会揭晓。

现在把问题集中起来：原子怎样排列，才能同时解释“规则外形”“固定熔点”“能剥成薄片”这三件事？

\section{把镜头推近：原子在排队}

把一粒食盐放大一亿倍，你会看到第 18 章讲过的景象：钠离子和氯离子像棋盘格一样交替排列，向上、向下、向前、向后、向左、向右，全都以完全相同的节奏重复下去——一块食盐里没有任何一个“\chem{NaCl} 分子”，只有一张无限延伸的离子阵列。

这种在三维空间不断重复的规则阵列，叫做\textbf{晶格}；阵列里最小的重复单元，叫做\textbf{晶胞}。晶胞之于晶格，就像墙纸上的一个完整图案单元之于整张墙纸：只要知道了晶胞长什么样、朝哪边重复，整块晶体的结构就被一条规则全部写完了。换个说法：晶体内部是\textbf{长程有序}的——只要测量出一个粒子的位置，你就能预言它第几亿个邻居在哪里。

\begin{center}
\begin{tikzpicture}[scale=0.85]
\foreach \x in {0,0.6,1.2,1.8,2.4}
  \foreach \y in {0,0.6,1.2,1.8}
    \fill[mainblue] (\x,\y) circle (0.1);
\draw[mainblue,dashed] (0.6,0.6) rectangle (1.2,1.2);
\node at (1.2,-0.6) {晶体：处处重复（虚线是一个晶胞）};
\begin{scope}[xshift=5.4cm]
\fill[warnred] (0.15,0.25) circle (0.1); \fill[warnred] (0.65,0.55) circle (0.1); \fill[warnred] (0.35,1.15) circle (0.1);
\fill[warnred] (1.05,0.05) circle (0.1); \fill[warnred] (1.25,0.85) circle (0.1); \fill[warnred] (0.85,1.55) circle (0.1);
\fill[warnred] (1.65,0.45) circle (0.1); \fill[warnred] (1.55,1.35) circle (0.1); \fill[warnred] (2.15,0.95) circle (0.1);
\fill[warnred] (2.05,1.75) circle (0.1); \fill[warnred] (0.2,1.8) circle (0.1); \fill[warnred] (2.25,0.15) circle (0.1);
\node at (1.2,-0.6) {玻璃：挤在一起，却没有队列};
\end{scope}
\end{tikzpicture}
\end{center}
（示意图：圆点是人为的表示法，不代表原子的真实外观和颜色。）

晶体的三大怪癖，用排队模型一口气全解释通了：

\begin{itemize}[itemsep=2pt]
\item \textbf{规则外形}：食盐的晶胞是立方体，沿三个互相垂直的方向整整齐齐地堆，堆出来的宏观颗粒自然还是立方体。晶面就是排列最整齐的那些“切面”。
\item \textbf{能剥成薄片}：云母层内粒子结合得牢，层与层之间结合得松，受力时当然从最弱的地方——层间——断开。
\item \textbf{固定熔点}：阵列里每个粒子的处境一模一样，邻居一样多、距离一样远，所以束缚一样紧。加热到某个温度，大家几乎同时“解锁”，熔点自然是一条清晰的线。
\end{itemize}

而玻璃里的粒子是右边那幅图的样子：它们也紧紧挤在一起（密度和晶体差不了太多），但谁和谁挨着完全没有章法——这就是\textbf{长程无序}。

\begin{qiaoliang}
一粒盐里有多少个晶胞？来算一笔账。食盐的晶胞是边长约 \ud{0.56}{nm} 的立方体。取一粒边长约 \ud{0.5}{mm} 的盐粒，每条边上能排的晶胞数是
\[
\frac{0.5\times 10^{-3}}{0.56\times 10^{-9}} \approx 9\times 10^{5}\ \text{（约九十万个）}
\]
立方之后，整粒盐约有 $(9\times 10^{5})^{3} \approx 7\times 10^{17}$ 个晶胞。每个晶胞含 4 个钠离子和 4 个氯离子，所以这粒肉眼几乎看不清的盐里，大约有 $6\times 10^{18}$ 个离子——六十亿亿个。

再看“有序”的真正含义：这 $10^{18}$ 个粒子的位置，用“一个晶胞 $+$ 三个重复方向”这一条规则就能全部写出来，信息量小到可以印在邮票上。而同样多粒子组成的一块玻璃，每个粒子的位置都得单独记录——有序与无序的差别，首先是信息量的差别。
\end{qiaoliang}

\section{为什么排队？又为什么来不及排队}

第 17 章说过一个朴素的原则：当变化能让整个体系的总能量降低时，它就有可能发生——原子没有愿望，驱动一切的是能量和概率。粒子排成晶格时，每个粒子吸引的邻居最多、距离最合适、排斥最小，体系的总能量最低。所以从能量方向上看，大多数纯净物质冷却下来都“该”结晶。

但排队需要时间。液体里的粒子在不停游荡，要刚好撞到正确的位置、转到正确的角度，才能嵌进阵列。这里就出现了全书会反复登场的一对分工——\textbf{方向与速度}：

\begin{itemize}[itemsep=2pt]
\item \textbf{慢慢冷却}：粒子有充足时间各就各位，能长出大块晶体。地下的岩浆用几万年慢慢冷凝，花岗岩里就长出了肉眼可见的石英和长石晶体。
\item \textbf{急速冷却}：粒子还没归位就被“冻”在原地，形成玻璃。岩浆喷进冷水，几分钟内凝固，就成了黑曜石——一种天然玻璃。
\end{itemize}

换句话说：结晶是能量上更合算的\textbf{方向}，玻璃是来不及走到终点就被定在半路的\textbf{产物}。这不是“放热就自动发生”的反例，而是热力学管方向、动力学管速度的教科书式案例——到第 27 章和第 29 章，这对分工还会回来，而且会成为主角。

巧克力师傅每天都在实践这个原理。可可脂能排出好几种晶型，只有一种能给出脆响、光泽和入口即化的口感。所谓“调温”，就是精确地升温、降温、再回温，逼脂肪分子排成那一种队列。而巧克力放久后的白霜，是脂肪在常温下用几个月时间慢慢重排成了更稳定的队形——能量上更合算，口感上更难吃。

\section{四种排队方式}

粒子排队时，用什么“胶水”粘住彼此，决定了晶体的脾气。按胶水不同，晶体分成四大类：

\begin{table}[htbp]
\centering
\begin{tabular}{lllll}
\toprule
类型 & 排队的粒子 & “胶水” & 例子 & 脾气 \\
\midrule
离子晶体 & 正、负离子 & 静电吸引 & \chem{NaCl}、\chem{CaF_2} & 硬而脆，熔点高 \\
分子晶体 & 完整分子 & 分子间作用力 & 冰、干冰、糖 & 软，熔点低 \\
共价网络晶体 & 原子 & 共价键织网 & 金刚石、石英 & 极硬，熔点极高 \\
金属晶体 & 金属阳离子 & 电子海 & \chem{Cu}、\chem{Fe} & 导电，延展 \\
\bottomrule
\end{tabular}
\end{table}

分子晶体的“胶水”是第 22 章讲的分子间作用力——比化学键弱一两个数量级，所以干冰（固态 \chem{CO_2}）在 $-78\,^{\circ}\mathrm{C}$ 就直接升华，冰在第 23 章讲过靠氢键支撑也只撑到 $0\,^{\circ}\mathrm{C}$。金属晶体的“电子海”是第 21 章的主角，导电、导热、能拉成丝，全靠它。离子晶体在第 18 章已经详细逛过，这里只提醒一句：它的硬来自处处有吸引，它的脆来自一旦错动半格、同性电荷对撞就全盘崩裂。

真正的新面孔是\textbf{共价网络晶体}：整个晶体就是一张共价键织成的三维大网，一块石英可以说就是一个“分子”——一个分子量大到天文数字的巨分子。要熔化它，没有“先松开哪里”的余地，只能大量剪断共价键（断键要吸能，第 17 章），所以石英要烧到约 $1700\,^{\circ}\mathrm{C}$ 才熔化。

这里符号终于可以登场了，而且一登场就要纠一个误会：石英写作 \chem{SiO_2}，但世界上并不存在一个独立的“二氧化硅分子”。石英里每个硅原子连着 4 个氧，每个氧连着 2 个硅，向四面八方无限延伸；化学式 \chem{SiO_2} 只记录硅和氧的最简个数比 $1:2$。这和 \chem{NaCl} 只记录 $1:1$ 是同一个道理——对网络晶体和离子晶体，化学式是“配料比例”，不是“分子画像”。

\begin{shiyan}
\textbf{养一颗属于自己的晶体}。材料：明矾（食品添加剂店或网店有售，也可用粗盐代替）、一个干净玻璃杯、热水、一根棉线、一根筷子、一颗形状规整的小晶粒（作“晶种”）。步骤：①向热水中分批加入明矾并搅拌，直到再有少量不再溶解，得到接近饱和的溶液；②把溶液倒入玻璃杯静置，第二天杯底会结出小晶体，挑出最规整的一颗当晶种；③把晶种系在棉线上，棉线另一头缠在筷子上，让晶种悬在溶液中央，杯口盖张纸防尘；④放在不受震动的角落，等一到两周。观察点：晶体一天天长出平整的新晶面，形状始终保持着对称的样子；对比棉线上偶然结出的杂乱小晶粒。想一想：为什么晶种能“指挥”后来的粒子按同样的队形排列？安全提示：热水小心烫伤；明矾溶液和晶体都\textbf{不能吃}，做完实验洗手；废液倒入下水道即可。
\end{shiyan}

\section{同样是碳：金刚石与石墨}

现在可以正面回答开场的第一问了。金刚石和石墨都是纯碳，差别只在原子怎么排队。

\textbf{金刚石}里，每个碳原子伸出 4 只“手”（回忆第 20 章的四面体构型），拉住周围 4 个邻居，邻居再拉邻居的邻居，向上下左右前后均匀延伸成一张刚性的三维网。这张网没有薄弱环节：想让原子滑动，往哪个方向推都要同时剪断许多共价键。所以金刚石是莫氏硬度表上的满分选手，能划开玻璃和石头。而且每个碳的 4 个电子全部锁在共价键里上班，没有一个自由电子——所以纯净的金刚石不导电，还透明：可见光的能量不足以撼动那些被锁死的电子，光直接穿过去了。它折光的本领又极强，工匠把钻石切成几十个精确角度的小面，让射进去的光在内部反复全反射后再射出，钻石才闪得那么耀眼。顺带一提，金刚石还是自然界导热最好的材料之一，摸上去总有股凉意——珠宝店用来鉴别真钻的“热导仪”，测的就是这一点。

\textbf{石墨}里，每个碳只连 3 个邻居，织成一张张六边形蜂窝平面。层内是极强的共价键（比金刚石里的还短、还牢），层与层之间却只靠第 22 章那种微弱的分子间作用力维系，层间距 \ud{0.335}{nm}。于是层与层一推就滑——这就是石墨软、能写字、能做铅笔芯和润滑剂的全部原因。同时，每个碳多出 1 个没在层内成键的电子，这些电子在层内自由跑动，所以石墨沿层的方向能导电。

铅笔写字的原理于是真相大白：不是铅笔芯被“磨掉”，而是纸面像铲子一样，把石墨的蜂窝层一层一层刮下来留在纸上。一道笔画就算只有 \ud{100}{nm} 厚，也叠着大约三百层蜂窝片——你随手一划，就在纸上叠了一座三百层的“分子楼”。

顺便回答第二问：铅笔芯里\textbf{没有铅}。它是石墨粉加黏土烧制出来的，黏土比例高就硬（标号 H），石墨比例高就黑（标号 B）。古人曾把石墨误认成铅矿石，“铅笔”这个将错就错的名字一直用到今天。

而把石墨撕到只剩\textbf{一层}，就得到一种名声大噪的新材料——石墨烯。它的故事先卖个关子，放到新材料那一节讲。

\section{玻璃：被速冻的乱局，但不是缓慢流动的液体}

普通玻璃的主要成分也是 \chem{SiO_2}（加碳酸钠和石灰石降低熔点、改善性能），和石英的区别只在排列。把石英玻璃的镜头推进去看：局部其实很有秩序——每个硅原子照样被 4 个氧围成四面体，邻居关系规规矩矩；但四面体与四面体连成了大大小小、歪歪扭扭的环，整体找不到任何重复周期。这叫\textbf{短程有序、长程无序}。

这种结构立刻解释了玻璃的全部怪脾气：

\begin{itemize}[itemsep=2pt]
\item \textbf{没有固定熔点}：阵列里每个粒子的处境各不相同，束缚有松有紧。加热时松的先动、紧的后动，于是一路连续软化——玻璃工匠正是趁这团“软糖”状态吹出瓶子和灯罩。
\item \textbf{各方向性质相同}：没有整齐的弱面可循，敲碎时断面只能随机弯曲，形成贝壳状纹路。
\item \textbf{它是如假包换的固体}：常温下玻璃有固定形状，受力能弹回，受猛力会碎——固体的力学行为它一样不缺。准确的叫法是\textbf{无定形固体}：结构像液体，力学是固体。
\end{itemize}

\begin{wuqu}
“玻璃其实是流动极慢的液体，老教堂窗户下厚上薄就是证据。”——这是个流传极广的误会，三条反例足以拆穿它。第一，流动量算不出来：玻璃的黏度随温度下降指数式暴涨，在室温下已大到流动完全停滞，就算等上远超宇宙年龄的时间，也流不出肉眼可见的变形。第二，精密仪器作证：两三百年前的望远镜和显微镜镜片，磨好之后至今仍精确成像；两千年前的古罗马玻璃瓶，也没见哪个“流”成了上薄下厚。第三，厚薄不匀另有原因：老式窗玻璃用吹制法或转盘法制造，天生一边厚一边薄，工匠装窗时习惯把又厚又重的边朝下——那是安装习惯，不是流动的下场。诚实地说：玻璃是“被冻住的液体结构”，不是“还在流动的液体”。
\end{wuqu}

\section{科学家怎样知道原子在排队}

晶胞的边长只有 \ud{0.56}{nm}，比可见光的波长还小几百倍，任何显微镜都看不见。那我们凭什么确信原子真的排成了方阵？

\begin{zhengju}
1912 年，德国物理学家劳厄想到一个漂亮的检验。当时 X 射线刚被发现十几年，没人确定它是什么。劳厄推理：\textbf{如果}晶体真是粒子的规则阵列，\textbf{而} X 射线是波长与原子间距同数量级的波，那么晶体对 X 射线就该像一块三维的光栅，把 X 射线衍射成规则的斑点图案。他的助手弗里德里希和尼平拿硫酸铜晶体做了实验——底片上果然出现排列整齐的斑点。一个实验同时证明了两件事：X 射线是波，晶体内部有周期排列。要知道，当时“原子真实存在”都还被一些科学家怀疑。

布拉格父子接着把斑点变成了尺子。儿子劳伦斯·布拉格发现，可以把衍射看成原子平面层对 X 射线的“反射”，由斑点的位置就能反推出原子平面之间的间距。父子俩解出的第一批结构里就有食盐 \chem{NaCl}——结果出人意料：食盐里根本找不到一个个成对的“\chem{NaCl} 分子”，只有正负离子交替的无限阵列。这正是第 18 章那句“离子晶体是巨大排列”的实验来源。1915 年，父子俩共同获得诺贝尔奖，儿子当时只有 25 岁。

这套方法的寿命比任何单个结论都长：X 射线衍射后来被用来解读蛋白质、DNA 等生命大分子的结构——这本书讲到遗传的章节，你还会和这门手艺重逢。这个故事最该记住的是方法：科学家没有“看见”原子，而是设计了一种原子必然留下指纹的局面，再从指纹反推现场。
\end{zhengju}

\section{新材料：换成分，不如换结构}

二十世纪材料科学最大的教训可以写成一句话：\textbf{性质不只由“含什么元素”决定，更由结构在各个尺度上怎样组织决定。}

\begin{itemize}[itemsep=2pt]
\item \textbf{石墨烯}：单层石墨。2004 年，盖姆和诺沃肖洛夫用一卷透明胶带反复粘贴、撕开石墨，硬是把蜂窝层撕到只剩一层，2010 年因此获得诺贝尔物理学奖。一层碳原子的薄片，同样重量下强度可达钢的上百倍，导电导热性能极佳。眼下的难题不在“它好不好”，而在“怎样便宜地造出大片完美的单层”——又一个结构与速率的问题。
\item \textbf{陶瓷}：氧化铝 \chem{Al_2O_3}、氮化硅 \chem{Si_3N_4} 这类材料，靠离子键与共价键织成网络，硬、耐高温、耐腐蚀、不导电，于是有了陶瓷刀、火花塞和航天器的隔热瓦。代价是脆：网络键不允许粒子滑移，裂纹一出现就一传到底。陶瓷与金属的性格差异，几乎全部是“键的种类决定粒子能否滑移”这一个差异的推论。
\item \textbf{复合材料}：单一材料总有短板，那就拼装。钢筋抗拉、混凝土抗压，捆在一起成了钢筋混凝土；玻璃纤维嵌进树脂成了玻璃钢；碳纤维嵌进树脂，造出比铝轻、比钢强的自行车架和球拍。复合材料的性能不只写在成分表里，更写在“谁站在哪个位置、朝哪个方向”里。
\item \textbf{超材料}：这条路被推到极端——性质主要来自人工设计的微结构，而不是成分。把比波长还小的结构单元按特定图案排列，材料对电磁波、声波的响应能出现自然界没有的脾气，比如让光“拐弯绕开”一块区域的初步隐身装置。在这里，成分几乎退场，几何登场。
\end{itemize}

这些新材料早已在你身边服役：芯片是从一整根超纯单晶硅棒上切下来的——整根棒是一个完整晶格，一粒杂质原子都能改变它的导电脾气；人工髋关节的球头用氧化锆陶瓷，看中的是它的硬和耐磨；药厂用 X 射线衍射看清致病蛋白质的形状，再设计刚好能卡进去的药物分子。从原子排列（纳米尺度）到晶粒组织（微米尺度）再到零件外形（宏观尺度），\textbf{结构跨着尺度决定材料的一切}——这就是本章最后要钉进你脑子的那句话。

\section{边界：完美的晶体并不存在}

本章讲了半天的“完美晶格”，其实是个理想模型。真实晶体满是缺陷：少一个粒子的“空位”、硬挤进来的杂质原子、整排错开半格的“位错”。

但缺陷不等于瑕疵，很多重要的性质恰恰来自缺陷：金属能被锻打弯曲，靠的就是位错一行一行地滑移；铝合金、钢比纯金属硬，是因为杂质原子把位错“卡”住了；半导体的灵魂更是故意掺进去的杂质——没有缺陷就没有芯片。

还有两件事突破了本章的框架。其一，你手里的铁钉不是一个大晶体，而是无数取向各异的小晶粒拼成的\textbf{多晶}，晶粒之间的边界深刻影响着强度和生锈的方式。其二，1982 年，谢赫特曼在一种铝合金的衍射图案里看到了五次对称——教科书当时断言“周期排列不可能有五次对称”，他因此被同行嘲笑了多年。后来这种\textbf{准晶体}在天然陨石里也被找到，2011 年他拿到了诺贝尔化学奖。有序，原来不必等于周期重复。所以请把“完美晶格”当作思维脚手架：它很好用，但真实材料的故事，有一半写在缺陷和边界里。

\section{回到那支铅笔}

兑现承诺的时候到了。

铅笔能写字，因为石墨层内是强共价键、层间只有微弱的分子间作用力，纸面一刮，蜂窝层成片剥落；钻石最硬，因为每个碳都被四面体共价网从四面八方锁死，无处可滑。铅笔芯里没有铅，只有石墨和黏土。你家那扇玻璃窗，是一锅被速冻的 \chem{SiO_2} 乱局——短程有序、长程无序，常温下它是如假包换的固体；老教堂窗户下厚上薄，记的是工艺和安装习惯的账，不是流动的账。

同一种元素、两种极端性格——现在你明白了，分水岭从来不在成分表上，而在三个问题上：粒子怎样排列、用什么键连接、组织到哪个尺度。这也是那五个贯穿问题里前两个的又一次实战。

\section{下一章：把晶格拆开}

晶格把粒子锁得整整齐齐，能量上还非常合算。可一把盐撒进水里，几秒钟就“消失”了。水分子既不带电锤也不带撬棍，凭什么拆掉一座离子方阵？拆走的离子去了哪里、换上了什么装束？为什么有的晶体（食盐）一冲就散，有的（沙子里的石英）泡一万年也纹丝不动？

第 25 章《溶液里发生了什么》，我们跟着一颗钠离子下水。

\begin{tiaozhan}
布拉格父子的“尺子”其实只有一行公式。把晶体里间距为 $d$ 的一层层原子平面看成半透明的镜子，X 射线以角度 $\theta$ 射入，相邻两层反射波的波程差恰好是 $2d\sin\theta$。当波程差等于波长的整数倍时，反射波互相加强，底片上出现亮斑：
\[
n\lambda = 2d\sin\theta
\]
这就是布拉格定律。动手算一算：实验常用的 X 射线波长 $\lambda = \ud{0.154}{nm}$，对食盐测得一级亮斑（$n=1$）出现在 $\theta \approx 16^{\circ}$，那么对应的原子面间距 $d = \frac{\lambda}{2\sin\theta} \approx \frac{0.154}{2\times 0.276} \approx \ud{0.28}{nm}$——正好是晶胞边长 \ud{0.56}{nm} 的一半，因为那组原子面把晶胞对半切开。一个斑点，换算出一排原子的位置；几百个斑点，就能拼出整块晶体的三维结构。跳过本节不影响主线，但你会错过“尺子”最精妙的部分。
\end{tiaozhan}

\section*{练一练}

\begin{sikaoti}
\item 画三幅粒子示意图（用圆圈代表粒子），分别表现：单晶、玻璃、多晶（几个取向不同的小晶粒拼在一起）。每幅图下写一句判据：别人凭什么看出它是哪一种？
\item 干冰（分子晶体）和石英（共价网络晶体）都是非金属氧化物的固体。从“排队方式和胶水种类”出发，解释为什么干冰在 $-78\,^{\circ}\mathrm{C}$ 就升华，石英却要到约 $1700\,^{\circ}\mathrm{C}$ 才熔化。
\item 石墨烯和金刚石都只含碳原子。从“电子在哪里上班”的角度回答：为什么一个导电、一个不导电？
\item 有人坚持要把玻璃归为液体。列出本章给出的两条反驳证据，并给出你认为更准确的分类名称和理由。
\item 用“脂肪分子换了排队方式”解释：起白霜的巧克力为什么是安全的、口感却变沙了？想要恢复脆亮的口感，巧克力师傅该控制什么变量？
\item 估算：食盐晶胞边长约 \ud{0.56}{nm}，每个晶胞含 4 个钠离子。一粒边长约 \ud{1}{mm} 的立方体食盐里，大约有多少个钠离子？写出你的计算过程。
\end{sikaoti}
